% Generated draft: requires mathematical and source review.
\documentclass[11pt]{article}
\usepackage[margin=25mm]{geometry}
\usepackage{fontspec}
\setmainfont{Latin Modern Roman}
\usepackage{amsmath,amssymb,mathtools,unicode-math}
\setmathfont{Latin Modern Math}
\usepackage{enumitem,needspace,xurl,hyperref}
\setlength{\parindent}{0pt}
\setlength{\parskip}{5pt}
\setlength{\emergencystretch}{3em}
\newcommand{\N}{\mathbb{N}}
\newcommand{\Z}{\mathbb{Z}}
\newcommand{\Q}{\mathbb{Q}}
\newcommand{\R}{\mathbb{R}}
\newcommand{\C}{\mathbb{C}}
\newcommand{\F}{\mathbb{F}}
\newcommand{\A}{\mathbb{A}}
\newcommand{\PP}{\mathbb{P}}
\newcommand{\E}{\mathbb{E}}
\newcommand{\eps}{\varepsilon}
\newcommand{\dd}{\,\mathrm{d}}
\begin{document}
\section*{UnsolvedMath ID: 30001969}
\textbf{Source problem code:} catalogue entry
\subsection{Definitions and mathematical statement}
\textbf{Supplied discussion question.} ``How can algorithms in medieval mathematical texts that are obscure to modern readers be understood, and what elements guarantee their correctness?''

The supplied record concerns interpreting medieval mathematical procedures and identifying grounds for their correctness. It specifies no particular text or algorithm, input domain, admissibility conditions, intended output, executable interpretation, or criterion of correctness. Consequently, no definite quantified mathematical proposition or correctness assertion can be recovered without additional source evidence. No algorithm or correctness theorem is asserted here.

\textbf{Unresolved qualification.} The supplied statement status is \texttt{unrecoverable}. Its statement-verification note characterizes the workshop material as a historical, pedagogical, or methodological discussion prompt rather than an unresolved mathematical proposition; no separate clean statement is supplied. The registry nevertheless labels the record \texttt{open}. That label is supplied curation, not verification that a specific mathematical conjecture remains open.
\subsection{Short English statement}
The question asks how to interpret obscure medieval mathematical algorithms and establish their correctness. The record supplies no specific algorithm or correctness requirement, so a precise mathematical claim cannot be extracted. This statement-only draft requires human review and further source evidence.
\subsection{Sources}
\begin{itemize}
\item Local UnsolvedMath registry record, permanent ID \(30001969\), exact source problem code \texttt{catalogue entry}, title \emph{Interpreting Obscure Medieval Mathematical Algorithms}; collection: Oberwolfach Reports Open Problems. The supplied proposed year is 2012; no individual proposer is identified.
\item Supplied source citation: \emph{Explicit Versus Tacit Knowledge in Mathematics} (2012), Oberwolfach Reports, pp.\ 131--210, DOI \texttt{10.4171/owr/2012/04}. Supplied source URL and literature-source link, labeled ``original Oberwolfach report'': \url{https://doi.org/10.4171/owr/2012/04}. The record attributes the question to workshop open-problem discussions documented in this report; no item-specific page or author is supplied.
\item Literature assessment dated \texttt{2026-08-21}: the supplied curation labels the record open and states that the questions concern interpretation and correctness of medieval algorithms, with no single formal claim in the extraction whose resolution could be verified. This is supplied dated curation derived from the Oberwolfach Reports record, not an independently verified literature assessment. The supplied background requests independent verification of time-sensitive status and cited claims.
\end{itemize}
\end{document}
